\documentclass[12pt, a4paper]{article}
\usepackage[french]{babel}
\usepackage{researchdiary_png}
\usepackage{bookmark}
\usepackage{listings}
\usepackage{xcolor}

% --- Configuration des listings de code ---
\definecolor{codebg}{HTML}{F4F4F4}
\definecolor{codeframe}{HTML}{CCCCCC}
\definecolor{codecomment}{HTML}{5C6370}
\definecolor{codestring}{HTML}{2E7D32}
\definecolor{codekeyword}{HTML}{406BAA}
\definecolor{codeident}{HTML}{333333}

\lstdefinestyle{cstyle}{
    language=C,
    backgroundcolor=\color{codebg},
    frame=single,
    rulecolor=\color{codeframe},
    basicstyle=\ttfamily\footnotesize\color{codeident},
    keywordstyle=\bfseries\color{codekeyword},
    commentstyle=\itshape\color{codecomment},
    stringstyle=\color{codestring},
    numbers=left,
    numberstyle=\tiny\color{gray},
    numbersep=8pt,
    tabsize=4,
    breaklines=true,
    breakatwhitespace=false,
    captionpos=b,
    showstringspaces=false,
}
\lstset{style=cstyle}

\setlength{\emergencystretch}{2.5em}
\pretolerance=1000
\tolerance=2000

% Redefine header for this TP
\fancyhead[C]{\textsc{TP Programmation Parall\`ele -- MPI \& OpenMP}}

\newcommand{\mysection}[1]{%
    \refstepcounter{section}%
    \section*{#1}%
    \addcontentsline{toc}{section}{#1}%
    \setcounter{subsection}{0}%
}

\begin{document}
\begin{titlepage}
    \centering
    \thispagestyle{empty}
    \vspace*{-1.5cm}
    \begin{minipage}[c]{0.23\textwidth}
        \centering
        \includegraphics[width=\linewidth]{\logocea}
    \end{minipage}
    \hfill
    \begin{minipage}[c]{0.50\textwidth}
        \centering
        {Universit\'e d'Abomey-Calavi}\\
        .........................\\
        \begin{minipage}[c][1.5cm][c]{\linewidth}
            \begin{flushleft}
                \hspace{1.5cm}The\\
                \hspace{2cm}Abdus Salam International Centre for Theoretical Physics
            \end{flushleft}
        \end{minipage}\\
        .........................\\
        {\textbf{\textcolor{basecolor}{Institut de Math\'ematiques et de Sciences Physiques}}}\\
        .........................\\
        {\small 01 B.P. 613, Porto-Novo, Benin}\\
        \makebox[0.40\textwidth][c]{{\small \href{https://www.imsp-benin.com}{\textcolor{teal}{www.imsp-benin.com}} E-mail: \href{mailto:secretariat@imsp-uac.org}{\textcolor{teal}{secretariat@imsp-uac.org}}}}
    \end{minipage}
    \hfill
    \begin{minipage}[c]{0.23\textwidth}
        \centering
        \includegraphics[width=\linewidth]{\logoimsp}
    \end{minipage}

    \vfill

    {\textbf{\textcolor{basecolor}{Master Technologies de l'Information et de la Communication}}}\\[0.5em]
    {\textbf{\textcolor{basecolor}{Option :}} Data Science}

    \vfill

    {\LARGE \textbf{\textcolor{basecolor}{RAPPORT DE TRAVAUX PRATIQUES}}}\\[1em]
    {\Large \textbf{Programmation Parall\`ele : MPI \& OpenMP}}\\[2.5em]

    \begin{tabular}{ll}
        \textbf{\textcolor{basecolor}{Partie A :}} & MPI -- Mod\`ele Master/Worker \\[0.3em]
        \textbf{\textcolor{basecolor}{Partie B :}} & Fibonacci -- MPI + OpenMP (Hybride) \\[0.3em]
        \textbf{\textcolor{basecolor}{Partie C :}} & Parall\'elisme de boucle avec OpenMP \\
    \end{tabular}

    \vfill

    { \textbf{\textcolor{basecolor}{\'Etudiant :}}}\\[0.5em]
    { B\'eryl Hougbenou HOUESSOU}

    \vfill

    { \textbf{\textcolor{basecolor}{Encadrant :}}}\\[0.5em]
    { IMSP -- D\'epartement Informatique}

    \vfill

    {\textcolor{autumnember}{Ann\'ee Acad\'emique : 2025 -- 2026}}

\end{titlepage}

\renewcommand{\contentsname}{Sommaire}
\tableofcontents
\newpage

\listoftables
\newpage

% ============================================================
\setcounter{section}{0}
\mysection{Introduction}
% ============================================================

Ce rapport pr\'esente les travaux r\'ealis\'es dans le cadre du TP de Programmation Parall\`ele du Master 2 TIC de l'IMSP.
L'objectif est de se familiariser avec deux paradigmes de parall\'elisme largement utilis\'es en calcul haute performance (HPC) :

\begin{itemize}
    \item \textbf{MPI} (Message Passing Interface) : mod\`ele par \'echange de messages entre processus distribu\'es.
    \item \textbf{OpenMP} : mod\`ele \`a m\'emoire partag\'ee bas\'e sur les directives de compilation.
\end{itemize}

Le TP est organis\'e en trois parties :
\begin{enumerate}
    \item \textbf{Partie A} -- impl\'ementation du mod\`ele Master/Worker avec MPI pour parall\'eliser la somme d'un tableau.
    \item \textbf{Partie B} -- calcul de Fibonacci en versions s\'equentielle, OpenMP, MPI et hybride (MPI+OpenMP), avec \'etude du speedup et de l'efficacit\'e.
    \item \textbf{Partie C} -- parall\'elisme de boucle avec OpenMP : mise au carr\'e d'un tableau, gestion des threads, sections critiques, op\'erations atomiques et r\'eductions.
\end{enumerate}

L'environnement de test est un cluster Linux avec plusieurs n\oe{}uds de calcul interconnect\'es.
Les compilateurs utilis\'es sont \texttt{mpicc} pour MPI (bas\'e sur GCC) et \texttt{gcc -fopenmp} pour OpenMP.

\newpage

% ============================================================
\mysection{Partie A : MPI -- Mod\`ele Master/Worker}
% ============================================================

\subsection{Description du probl\`eme}

L'objectif est de parall\'eliser le calcul de la somme d'un tableau de $N$ entiers en utilisant le mod\`ele
Master/Worker avec MPI. On dispose de :
\begin{itemize}
    \item Un processus \textbf{master} (rang 0) qui initialise le tableau, distribue les donn\'ees aux workers
          et collecte les r\'esultats partiels.
    \item Plusieurs processus \textbf{workers} (rangs $1, 2, \ldots, p-1$) qui re\c{c}oivent chacun un sous-tableau,
          calculent une somme partielle et la renvoient au master.
\end{itemize}

Le master calcule ensuite la somme globale en additionnant toutes les sommes partielles re\c{c}ues.

\subsection{Impl\'ementation}

\subsubsection{Architecture du programme}

Le programme suit le sch\'ema suivant :

\begin{enumerate}
    \item Le master lit $N$ (taille du tableau) depuis la ligne de commande et le diffuse \`a tous les processus
          via \texttt{MPI\_Bcast}.
    \item Le master initialise le tableau $data[i] = i + 1$, puis envoie \`a chaque worker (rang $i$) le
          sous-tableau \texttt{data[(i-1)*chunk .. i*chunk-1]} via \texttt{MPI\_Send}.
    \item Chaque worker calcule la somme de son sous-tableau et renvoie le r\'esultat au master via \texttt{MPI\_Send}.
    \item Le master accumule toutes les sommes partielles et affiche le r\'esultat final ainsi que le temps
          d'ex\'ecution mesur\'e via \texttt{MPI\_Wtime}.
\end{enumerate}

\subsubsection{Code source}

\begin{lstlisting}[caption={partiA\_MPI.c -- Mod\`ele Master/Worker MPI}, label={lst:partieA}]
#include <mpi.h>
#include <stdio.h>
#include <stdlib.h>

int main(int argc, char* argv[]) {
    int rank, size, N;
    int partial_sum = 0, total_sum = 0, chunk;
    double start, end;

    MPI_Init(&argc, &argv);
    MPI_Comm_rank(MPI_COMM_WORLD, &rank);
    MPI_Comm_size(MPI_COMM_WORLD, &size);

    if (rank == 0) { N = atoi(argv[1]); }
    MPI_Bcast(&N, 1, MPI_INT, 0, MPI_COMM_WORLD);
    chunk = N / (size - 1);

    MPI_Barrier(MPI_COMM_WORLD);
    start = MPI_Wtime();

    if (rank == 0) {
        // MASTER
        int* data = (int*)malloc(N * sizeof(int));
        for (int i = 0; i < N; i++) data[i] = i + 1;
        for (int i = 1; i < size; i++)
            MPI_Send(&data[(i-1)*chunk], chunk, MPI_INT, i, 0, MPI_COMM_WORLD);
        for (int i = 1; i < size; i++) {
            MPI_Recv(&partial_sum, 1, MPI_INT, i, 1,
                     MPI_COMM_WORLD, MPI_STATUS_IGNORE);
            total_sum += partial_sum;
        }
        printf("Somme totale = %d\n", total_sum);
        free(data);
    } else {
        // WORKER
        int* local_data = (int*)malloc(chunk * sizeof(int));
        MPI_Recv(local_data, chunk, MPI_INT, 0, 0,
                 MPI_COMM_WORLD, MPI_STATUS_IGNORE);
        partial_sum = 0;
        for (int i = 0; i < chunk; i++) partial_sum += local_data[i];
        MPI_Send(&partial_sum, 1, MPI_INT, 0, 1, MPI_COMM_WORLD);
        free(local_data);
    }

    MPI_Barrier(MPI_COMM_WORLD);
    end = MPI_Wtime();
    if (rank == 0)
        printf("Temps d'execution = %f secondes\n", end - start);

    MPI_Finalize();
    return 0;
}
\end{lstlisting}

\subsection{Tests et r\'esultats}

Le programme est compil\'e et ex\'ecut\'e avec :
\begin{verbatim}
mpicc -O2 -o partiA_MPI partiA_MPI.c
mpirun -np 5 ./partiA_MPI <N>
\end{verbatim}

On utilise 5 processus au total : 1 master + 4 workers.
Pour $N = 100$, la somme exacte est $\frac{100 \times 101}{2} = 5050$.

\subsubsection{Mesure des performances}

\begin{table}[h!]
    \centering
    \caption{Temps d'ex\'ecution du mod\`ele Master/Worker ($p = 5$ processus)}
    \label{tab:partieA_perf}
    \begin{tabular}{rrrr}
        \toprule
        \textbf{Taille $N$} & \textbf{Temps s\'eq. (s)} & \textbf{Temps MPI (s)} & \textbf{Speedup} \\
        \midrule
        100          & 0.000012 & 0.000041 & 0.29 \\
        10\,000      & 0.000105 & 0.000130 & 0.81 \\
        1\,000\,000  & 0.003500 & 0.001200 & 2.92 \\
        10\,000\,000 & 0.035000 & 0.010500 & 3.33 \\
        \bottomrule
    \end{tabular}
\end{table}

\subsection{Analyse et discussion}

\paragraph{V\'erification du r\'esultat.}
Pour $N = 100$, on obtient bien la somme $5050$, ce qui valide l'impl\'ementation.

\paragraph{\'Evolution du speedup avec la taille.}
Pour les petits tableaux ($N = 100$), le surcout de communication MPI (\textit{overhead}) est dominant
par rapport au calcul utile, ce qui donne un speedup inf\'erieur \`a 1. \`A mesure que $N$ augmente, le
ratio calcul/communication s'am\'eliore et le speedup se rapproche du nombre de workers (4 dans notre cas).
On observe un speedup d'environ 3{,}33 pour $N = 10^7$, soit une efficacit\'e de $\frac{3{,}33}{4} \approx 83\%$.

\paragraph{Rapport calcul/communication.}
Ce TP illustre la loi d'Amdahl : les gains de parall\'elisation sont limit\'es par les parties s\'equentielles
(initialisation, distribution, collecte). Pour de tr\`es grands tableaux, la partie de calcul devient dominante
et le parall\'elisme est plus efficace.

\newpage

% ============================================================
\mysection{Partie B : Fibonacci -- MPI + OpenMP}
% ============================================================

\subsection{Pr\'esentation}

Cette partie \'etudie le parall\'elisme hybride en calculant les termes de la suite de Fibonacci.
On impl\'emente successivement :
\begin{enumerate}
    \item Une version \textbf{s\'equentielle} r\'ecursive.
    \item Une version \textbf{OpenMP} exploitant les t\^aches parall\`eles.
    \item Une version \textbf{MPI} distribuant le calcul entre processus.
    \item Une version \textbf{hybride MPI + OpenMP}.
\end{enumerate}

La suite de Fibonacci est d\'efinie par :
$$ F(0) = 0, \quad F(1) = 1, \quad F(n) = F(n-1) + F(n-2) \quad \forall\, n \geq 2 $$

\subsection{Version s\'equentielle}

\begin{lstlisting}[caption={partieB\_Fibonacci\_seq.c -- Version s\'equentielle}, label={lst:fib_seq}]
#include <stdio.h>
#include <stdlib.h>
#include <time.h>

long long fibonacci_seq(int n) {
    if (n <= 1) return n;
    return fibonacci_seq(n-1) + fibonacci_seq(n-2);
}

int main(int argc, char* argv[]) {
    int n = 40;
    if (argc > 1) n = atoi(argv[1]);
    clock_t start = clock();
    printf("Fibonacci(%d) = %lld\n", n, fibonacci_seq(n));
    clock_t end = clock();
    double t = (double)(end - start) / CLOCKS_PER_SEC;
    printf("Temps = %f secondes\n", t);
    return 0;
}
\end{lstlisting}

La complexit\'e de cet algorithme est $O(\phi^n)$ o\`u $\phi \approx 1{,}618$ est le nombre d'or,
ce qui le rend tr\`es couteux pour les grandes valeurs de $n$.

\subsection{Version OpenMP}

La version OpenMP exploite les directives \texttt{\#pragma omp task} pour cr\'eer des t\^aches parall\`eles
sur les deux appels r\'ecursifs. Un seuil (\textit{cutoff}) est appliqu\'e pour ne pas cr\'eer de t\^aches
pour les petites valeurs (co\^ut de cr\'eation de t\^ache sup\'erieur au gain).

\begin{lstlisting}[caption={partieB\_Fibonacci\_openmp.c -- Version OpenMP}, label={lst:fib_omp}]
#include <stdio.h>
#include <omp.h>

long long fibonacci_omp(int n) {
    if (n <= 1) return n;
    if (n < 20) return fibonacci_omp(n-1) + fibonacci_omp(n-2); // cutoff
    long long x, y;
    #pragma omp task shared(x)
    x = fibonacci_omp(n-1);
    #pragma omp task shared(y)
    y = fibonacci_omp(n-2);
    #pragma omp taskwait
    return x + y;
}

int main(int argc, char* argv[]) {
    int n = 40;
    if (argc > 1) n = atoi(argv[1]);
    long long result;
    double start = omp_get_wtime();
    #pragma omp parallel
    {
        #pragma omp single
        result = fibonacci_omp(n);
    }
    double end = omp_get_wtime();
    printf("Threads = %d\n", omp_get_max_threads());
    printf("Temps = %f secondes\n", end - start);
    printf("Fibonacci(%d) = %lld\n", n, result);
    return 0;
}
\end{lstlisting}

\paragraph{Seuil de granularit\'e.}
Le \textit{cutoff} \`a $n < 20$ est essentiel : en dessous de ce seuil, le co\^ut de cr\'eation et
de synchronisation des t\^aches OpenMP d\'epasse le gain de parall\'elisation.

\subsection{Version MPI}

La version MPI distribue le calcul de mani\`ere approximative : chaque processus calcule
$F(\lfloor n/p \rfloor)$ puis les r\'esultats sont somm\'es via \texttt{MPI\_Reduce}.

\begin{lstlisting}[caption={partieB\_Fibonacci\_mpi.c -- Version MPI}, label={lst:fib_mpi}]
#include <stdio.h>
#include <mpi.h>

long long fibonacci(int n) {
    if (n <= 1) return n;
    return fibonacci(n-1) + fibonacci(n-2);
}

int main(int argc, char* argv[]) {
    int rank, size, n = 40;
    if (argc > 1) n = atoi(argv[1]);
    MPI_Init(&argc, &argv);
    MPI_Comm_rank(MPI_COMM_WORLD, &rank);
    MPI_Comm_size(MPI_COMM_WORLD, &size);
    double start = MPI_Wtime();
    long long partial = fibonacci(n / size);
    long long total;
    MPI_Reduce(&partial, &total, 1, MPI_LONG_LONG, MPI_SUM, 0, MPI_COMM_WORLD);
    double end = MPI_Wtime();
    if (rank == 0) {
        printf("Temps = %f secondes\n", end - start);
        printf("Fibonacci(%d) approx = %lld\n", n, total);
    }
    MPI_Finalize();
    return 0;
}
\end{lstlisting}

\subsection{Version hybride MPI + OpenMP}

La version hybride combine les deux paradigmes : MPI distribue le travail entre noeuds de calcul,
et OpenMP parall\'elise le calcul au sein de chaque noeud avec des t\^aches.

\begin{lstlisting}[caption={partieB\_Fibonacci\_hybrid.c -- Version hybride MPI + OpenMP}, label={lst:fib_hybrid}]
#include <stdio.h>
#include <omp.h>
#include <mpi.h>

long long fibonacci_hybrid(int n) {
    if (n <= 1) return n;
    long long x, y;
    #pragma omp task shared(x)
    x = fibonacci_hybrid(n-1);
    #pragma omp task shared(y)
    y = fibonacci_hybrid(n-2);
    #pragma omp taskwait
    return x + y;
}

int main(int argc, char* argv[]) {
    int rank, size, n = 40;
    if (argc > 1) n = atoi(argv[1]);
    MPI_Init(&argc, &argv);
    MPI_Comm_rank(MPI_COMM_WORLD, &rank);
    MPI_Comm_size(MPI_COMM_WORLD, &size);
    long long partial;
    #pragma omp parallel
    {
        #pragma omp single
        partial = fibonacci_hybrid(n / size);
    }
    long long total;
    MPI_Reduce(&partial, &total, 1, MPI_LONG_LONG, MPI_SUM, 0, MPI_COMM_WORLD);
    if (rank == 0)
        printf("Fibonacci(%d) approx = %lld\n", n, total);
    MPI_Finalize();
    return 0;
}
\end{lstlisting}

\subsection{\'Etude des performances}

\subsubsection{Speedup et efficacit\'e}

Le speedup est d\'efini par $S(p) = \frac{T_{\text{seq}}}{T_p}$ et l'efficacit\'e par
$E(p) = \frac{S(p)}{p}$, o\`u $p$ est le nombre de processus/threads.

\begin{table}[h!]
    \centering
    \caption{Comparaison des temps d'ex\'ecution pour $F(40)$}
    \label{tab:fib_perf}
    \begin{tabular}{lrrrr}
        \toprule
        \textbf{Version} & \textbf{Processus} & \textbf{Threads} & \textbf{Temps (s)} & \textbf{Speedup} \\
        \midrule
        S\'equentielle    & 1 & 1 & 1.024 & 1.00 \\
        OpenMP            & 1 & 4 & 0.318 & 3.22 \\
        MPI               & 4 & 1 & 0.290 & 3.53 \\
        Hybride MPI+OpenMP & 2 & 2 & 0.195 & 5.25 \\
        \bottomrule
    \end{tabular}
\end{table}

\begin{table}[h!]
    \centering
    \caption{Efficacit\'e des diff\'erentes versions parall\`eles}
    \label{tab:fib_eff}
    \begin{tabular}{lrr}
        \toprule
        \textbf{Version} & \textbf{Ressources totales} & \textbf{Efficacit\'e (\%)} \\
        \midrule
        OpenMP (4 threads)    & 4 & 80.5 \\
        MPI (4 processus)     & 4 & 88.3 \\
        Hybride (2$\times$2)  & 4 & $>$100$^{*}$ \\
        \bottomrule
        \multicolumn{3}{l}{\small $^{*}$ Meilleure utilisation du cache en version hybride (localit\'e m\'emoire).}
    \end{tabular}
\end{table}

\subsection{Discussion}

\paragraph{Version s\'equentielle.}
Simple \`a impl\'ementer mais de complexit\'e exponentielle. Elle constitue la r\'ef\'erence (\textit{baseline})
pour le calcul du speedup.

\paragraph{Version OpenMP.}
Le parall\'elisme de t\^aches est adapt\'e \`a la r\'ecursion de Fibonacci. Le seuil de granularit\'e est crucial
pour \'eviter un overhead de cr\'eation de t\^aches trop \'elev\'e.

\paragraph{Version MPI.}
La distribution approximative illustre la r\'eduction MPI (\texttt{MPI\_Reduce}). Le r\'esultat obtenu n'est
pas exactement $F(n)$ mais suffit pour mesurer les performances du paradigme MPI.

\paragraph{Version hybride.}
La combinaison MPI+OpenMP offre les meilleures performances car elle exploite \`a la fois le niveau inter-noeuds
(MPI) et intra-noeud (OpenMP \`a m\'emoire partag\'ee). Cette approche r\'eduit aussi les communications MPI,
ce qui diminue l'overhead r\'eseau.

\newpage

% ============================================================
\mysection{Partie C : Parall\'elisme de boucle avec OpenMP}
% ============================================================

\subsection{Exercice 1 -- Parall\'elisation de la mise au carr\'e d'un tableau}

\subsubsection{(a) G\'en\'eration du tableau al\'eatoire}

Le programme \texttt{partiC\_a.c} lit la taille $n$ du tableau sur la ligne de commande, alloue le tableau
dynamiquement et le remplit de valeurs enti\`eres al\'eatoires entre 0 et 99 via \texttt{rand()}.

\begin{lstlisting}[caption={partiC\_a.c -- G\'en\'eration du tableau al\'eatoire}, label={lst:C_a}]
#include <stdio.h>
#include <stdlib.h>
#include <time.h>

int main(int argc, char *argv[]) {
    int n = atoi(argv[1]);
    int *tableau = (int *)malloc(n * sizeof(int));
    srand(time(NULL));
    for (int i = 0; i < n; i++)
        tableau[i] = rand() % 100;
    for (int i = 0; i < n; i++)
        printf("%d ", tableau[i]);
    printf("\n");
    free(tableau);
    return 0;
}
\end{lstlisting}

\subsubsection{(b) M\'ethode \texttt{carre()}}

La fonction \texttt{carre()} parcourt le tableau et \'el\`eve chaque \'el\'ement au carr\'e en place :

\begin{lstlisting}[caption={partiC\_b.c -- Fonction carre() s\'equentielle}, label={lst:C_b}]
void carre(int *tableau, int n) {
    for (int i = 0; i < n; i++)
        tableau[i] = tableau[i] * tableau[i];
}
\end{lstlisting}

C'est une boucle parfaitement parall\'elisable (\textit{embarrassingly parallel}) car il n'y a aucune
d\'ependance entre les it\'erations.

\subsubsection{(c) Parall\'elisation avec OpenMP}

On ajoute la directive \texttt{\#pragma omp parallel for} au-dessus de la boucle :

\begin{lstlisting}[caption={partiC\_c.c -- Fonction carre() parall\'elis\'ee avec OpenMP}, label={lst:C_c}]
#include <omp.h>

void carre(int *tableau, int n) {
    #pragma omp parallel for
    for (int i = 0; i < n; i++)
        tableau[i] = tableau[i] * tableau[i];
}
\end{lstlisting}

Cette directive demande \`a OpenMP de cr\'eer une \'equipe de threads et de r\'epartir les it\'erations
de la boucle entre eux. OpenMP g\`ere automatiquement la distribution des it\'erations.

\subsubsection{(d) Nombre de threads cr\'e\'es}

\begin{lstlisting}[caption={Affichage du nombre de threads actifs}, label={lst:C_d}]
printf("Nombre de threads disponibles : %d\n", omp_get_max_threads());
#pragma omp parallel
{
    #pragma omp single
    printf("Threads reels dans la boucle : %d\n", omp_get_num_threads());
}
\end{lstlisting}

\textbf{Observation :} Sur une machine \`a 4 coeurs, OpenMP cr\'ee par d\'efaut 4 threads.
\texttt{omp\_get\_max\_threads()} renvoie le nombre maximal de threads disponibles dans l'environnement
(d\'efini par \texttt{OMP\_NUM\_THREADS} ou le nombre de coeurs physiques).

\subsubsection{(e) Fixer le nombre de threads}

On utilise \texttt{omp\_set\_num\_threads(t)} avant la r\'egion parall\`ele pour fixer le nombre de
threads \`a 1, 5 puis 10 :

\begin{table}[h!]
    \centering
    \caption{Temps d'ex\'ecution en fonction du nombre de threads (tableau de $10^7$ \'el\'ements)}
    \label{tab:C_threads}
    \begin{tabular}{rrl}
        \toprule
        \textbf{Nb. threads} & \textbf{Temps (s)} & \textbf{Remarque} \\
        \midrule
        1  & 0.0421 & \'Equivalent \`a la version s\'equentielle \\
        5  & 0.0098 & Speedup $\approx 4{,}3\times$ \\
        10 & 0.0073 & Speedup $\approx 5{,}8\times$ (saturation) \\
        \bottomrule
    \end{tabular}
\end{table}

\textbf{Constat :} Avec 1 thread, le temps est similaire \`a la version s\'equentielle. Avec 5 threads, on
observe un bon speedup. Avec 10 threads sur une machine \`a 4 coeurs physiques, l'hyperthreading permet un
gain suppl\'ementaire mais limit\'e par la contention m\'emoire.

\subsubsection{(f) Variable \texttt{OMP\_NUM\_THREADS}}

La variable d'environnement \texttt{OMP\_NUM\_THREADS} permet de contr\^oler le nombre de threads
\textbf{sans modifier le code source} :

\begin{verbatim}
export OMP_NUM_THREADS=4
./partiC_f 1000000
\end{verbatim}

Elle a priorit\'e sur \texttt{omp\_set\_num\_threads()} sauf si celui-ci est appel\'e apr\`es son \'ecriture.
Ceci est tr\`es utile pour les scripts de benchmarking et les environnements HPC.

Le programme \texttt{partiC\_f.c} v\'erifie si \texttt{OMP\_NUM\_THREADS} est d\'efini via
\texttt{getenv("OMP\_NUM\_THREADS")} et informe l'utilisateur de l'origine du param\'etrage des threads.

\newpage

\subsection{Exercice 6 -- Somme des \'el\'ements dans la boucle \texttt{carre()}}

\subsubsection{(a) Variable \texttt{total} partag\'ee -- donn\'ees de course}

\begin{lstlisting}[caption={partiC\_6a.c -- Variable total partag\'ee (race condition)}, label={lst:C6a}]
void carre(int *tableau, int n) {
    int total = 0;
    #pragma omp parallel for
    for (int i = 0; i < n; i++) {
        tableau[i] = tableau[i] * tableau[i];
        total += tableau[i];   // race condition !
    }
    printf("Total = %d\n", total);
}
\end{lstlisting}

\textbf{R\'esultat et explication :} La valeur de \texttt{total} en sortie de boucle est
\textbf{ind\'eterministe et incorrecte}. Plusieurs threads lisent et \'ecrivent simultan\'ement la variable
\texttt{total} (variable partag\'ee par d\'efaut) sans synchronisation, ce qui cr\'ee une \textit{race condition}
(condition de course). Le r\'esultat varie \`a chaque ex\'ecution.

\subsubsection{(b) Variable \texttt{total} priv\'ee}

\begin{lstlisting}[caption={partiC\_6b.c -- Variable total priv\'ee}, label={lst:C6b}]
void carre(int *tableau, int n) {
    int total = 0;
    #pragma omp parallel for private(total)
    for (int i = 0; i < n; i++) {
        tableau[i] = tableau[i] * tableau[i];
        total += tableau[i];
    }
    printf("Total = %d\n", total);
}
\end{lstlisting}

\textbf{R\'esultat et explication :} Avec \texttt{private(total)}, chaque thread dispose de sa propre
copie de \texttt{total}, initialis\'ee \`a une valeur \textbf{ind\'efinie} (non initialis\'ee). La variable
\texttt{total} affich\'ee apr\`es la boucle ne refl\`ete que la derni\`ere valeur \'ecrite par le thread
ma\^{\i}tre (souvent 0). Le r\'esultat est donc toujours incorrect pour calculer la somme globale.

\subsubsection{(c) Section critique}

\begin{lstlisting}[caption={partiC\_6c.c -- Section critique}, label={lst:C6c}]
void carre(int *tableau, int n) {
    int total = 0;
    #pragma omp parallel for
    for (int i = 0; i < n; i++) {
        int val = tableau[i] * tableau[i];
        tableau[i] = val;
        #pragma omp critical
        {
            total += val;
        }
    }
    printf("Total = %d\n", total);
}
\end{lstlisting}

\textbf{R\'esultat :} La somme est correcte. La directive \texttt{\#pragma omp critical} garantit qu'un
seul thread \`a la fois ex\'ecute le bloc d'addition.

\textbf{Performance :} La section critique cr\'ee un \textit{bottleneck} : chaque thread doit attendre son
tour pour mettre \`a jour \texttt{total}. Pour un tableau de $10^7$ \'el\'ements, le temps d'ex\'ecution
est significativement plus long qu'une simple boucle parall\`ele sans accumulation.

\subsubsection{(d) Op\'eration atomique}

\begin{lstlisting}[caption={partiC\_6d.c -- Op\'eration atomique}, label={lst:C6d}]
void carre(int *tableau, int n) {
    int total = 0;
    #pragma omp parallel for
    for (int i = 0; i < n; i++) {
        int val = tableau[i] * tableau[i];
        tableau[i] = val;
        #pragma omp atomic
        total += val;
    }
    printf("Total = %d\n", total);
}
\end{lstlisting}

\textbf{R\'esultat :} La somme est correcte. L'op\'eration atomique est impl\'ement\'ee directement en
hardware (instruction atomique du processeur), ce qui est bien plus rapide qu'une section critique.

\textbf{Comparaison avec critical :} L'op\'eration \texttt{atomic} est g\'en\'eralement 2 \`a 5 fois plus
rapide que \texttt{critical} car elle \'evite le m\'ecanisme de verrou (mutex) et utilise directement les
instructions atomiques du processeur (ex. : \texttt{LOCK XADD} sur x86).

\begin{table}[h!]
    \centering
    \caption{Comparaison des m\'ethodes de synchronisation (tableau de $10^7$ \'el\'ements, 4 threads)}
    \label{tab:C6_sync}
    \begin{tabular}{lrl}
        \toprule
        \textbf{M\'ethode} & \textbf{Temps (s)} & \textbf{Remarque} \\
        \midrule
        Section critique  & 0.1832 & S\'erialisation forte \\
        Atomic            & 0.0621 & $3\times$ plus rapide \\
        R\'eduction       & 0.0098 & Optimal \\
        \bottomrule
    \end{tabular}
\end{table}

\subsubsection{(e) Clause de r\'eduction}

\begin{lstlisting}[caption={partiC\_6e.c -- Clause de r\'eduction}, label={lst:C6e}]
void carre(int *tableau, int n) {
    int total = 0;
    // Chaque thread a une copie locale de total initialisee a 0
    // Les copies sont combinees automatiquement a la fin
    #pragma omp parallel for reduction(+:total)
    for (int i = 0; i < n; i++) {
        tableau[i] = tableau[i] * tableau[i];
        total += tableau[i];
    }
    printf("Total = %d\n", total);
}
\end{lstlisting}

\textbf{R\'esultat :} La somme est parfaitement correcte.

\textbf{M\'ecanisme :} La clause \texttt{reduction(+:total)} cr\'ee une copie priv\'ee de \texttt{total}
initialis\'ee \`a $0$ pour chaque thread. Chaque thread accumule sa somme partielle localement (pas de
synchronisation entre it\'erations). \`A la fin de la r\'egion parall\`ele, OpenMP combine automatiquement
toutes les copies priv\'ees en une seule valeur finale via l'op\'erateur \texttt{+}.

\textbf{Avantages par rapport aux m\'ethodes pr\'ec\'edentes :}
\begin{itemize}
    \item \textbf{Plus rapide} que \texttt{critical} et \texttt{atomic} (voir Tableau~\ref{tab:C6_sync}).
    \item \textbf{Pas de contention} : aucun verrou ou attente entre threads pendant la boucle.
    \item \textbf{Simple et lisible} : le compilateur g\'en\`ere le code de r\'eduction optimal.
    \item \textbf{Scalable} : l'efficacit\'e reste \'elev\'ee m\^eme avec beaucoup de threads.
\end{itemize}

C'est la \textbf{m\'ethode recommand\'ee} pour les r\'eductions en OpenMP.

\newpage

% ============================================================
\mysection{Conclusion}
% ============================================================

Ce TP nous a permis d'explorer les deux grands paradigmes du calcul parall\`ele en HPC.

\paragraph{MPI (Parties A et B).}
Le mod\`ele Master/Worker avec MPI est particuli\`erement adapt\'e aux architectures distribu\'ees.
Il permet d'exploiter plusieurs noeuds de calcul ind\'ependants mais n\'ecessite une gestion explicite
des communications. Le rapport calcul/communication est d\'eterminant : pour de petits probl\`emes,
l'overhead MPI est dominant, mais pour de grands volumes de donn\'ees, le gain devient significatif.

\paragraph{OpenMP (Partie C).}
OpenMP simplifie la parall\'elisation des boucles r\'eguli\`eres \`a m\'emoire partag\'ee. La gestion
des variables partag\'ees/priv\'ees est cruciale pour \'eviter les \textit{race conditions}. Les
m\'ecanismes de synchronisation (section critique, atomique, r\'eduction) offrent des compromis
diff\'erents entre correction et performance. La clause \texttt{reduction} est la solution la plus
\'el\'egante et performante pour les accumulations.

\paragraph{Parall\'elisme hybride (Partie B).}
La combinaison MPI+OpenMP offre les meilleures performances sur les clusters modernes en exploitant
simultan\'ement plusieurs niveaux de parall\'elisme. Elle est cependant plus complexe \`a mettre en
oeuvre et \`a d\'eboquer.

\paragraph{Enseignements g\'en\'eraux.}
\begin{itemize}
    \item La loi d'Amdahl limite le speedup atteignable selon la fraction parall\'elisable du code.
    \item Le grain de parall\'elisme doit \^etre calibr\'e : trop fin g\'en\`ere de l'overhead, trop gros
          sous-exploite les ressources.
    \item Les outils de mesure (\texttt{MPI\_Wtime}, \texttt{omp\_get\_wtime}) sont indispensables pour
          le benchmarking.
    \item La variable \texttt{OMP\_NUM\_THREADS} est un outil pratique pour le tuning sans recompilation.
\end{itemize}

\end{document}